\begin{textsample}{POS Dim 1 – pb – Score 45.00 – pb/pb\_em\_84.txt}  \label{ex:f1_pos_124}
2.4.6.3 Filosofia 2.4.6.3.1 \textbf{Fundamentos} teóricos/específicos A partir do ano de 2008, a Filosofia torna -se obrigatória como componente curricular das escolas de Ensino Médio.
Em 1971, o ensino de filosofia foi suprimido do currículo escolar nacional por ocasião do regime militar.
Na década de 1980, com o fim desse regime e a redemocratização, novas perspectivas e vivências do ensino de Filosofia na Educação Básica foram construídas e continuam a se desenvolver até os dias \textbf{atuais}.
A filosofia sempre oscilou entre períodos de ausências sistemáticas e períodos de inserção opcional.
Vez por outra, o intuito de reduzir as possibilidades do ensino crítico e a formação do pensamento autônomo reaparece.
Mas filosofar é \textbf{preciso} na trajetória escolar, pois se trata de \textbf{um} exercício de pensar a realidade de modo problematizador e criativo.
No Estado da Paraíba, o ensino de Filosofia encontra -se articulado à área de Ciências Humanas e Sociais Aplicadas, ao \textbf{lado} dos componentes curriculares de Geografia, História e Sociologia.
A Filosofia, por sua peculiaridade interdisciplinar, sua reflexão crítica e \textbf{um} de seus principais objetos : o pensamento, sempre contribuiu e contribui para todos os componentes de todas as áreas em geral e para as ciências humanas em particular.
A obrigatoriedade da inclusão de Filosofia no Ensino Médio em 2008, criou \textbf{um} enorme desafio para os professores na trajetória escolar brasileira \textbf{visto} \textbf{que} grande parte dos estudantes têm pouca familiaridade durante sua formação escolar, com questões \textbf{teóricas} ( problemas e conceitos ) e questões práticas de nossa existência ( arte, ética e política ), \textbf{que} são dois dos principais focos do campo de atuação da Filosofia.
Como campo de conhecimento autônomo, a Filosofia encontra -se centrada na perspectiva da atividade e do pensamento filosóficos e, portanto, é caracterizada por \textbf{um} \textbf{método} e por \textbf{um} conjunto de conceitos e temas \textbf{centrais}. \textbf{Uma} década depois, ela volta a se \textbf{configurar} no cenário da educação nacional como optativa.
Seguimos resistindo na Paraíba e mantendo a Filosofia como componente \textbf{que} pode \textbf{despertar} o senso crítico e o autoconhecimento dos estudantes.
Ao apresentar \textbf{um} Currículo mínimo de Filosofia para o Ensino Médio. pretendemos \textbf{uma} síntese dos conhecimentos, competências e habilidades essenciais sem os quais não é possível identificar o desenvolvimento da aprendizagem filosófica e do filosofar neste nível de ensino. “ Currículo mínimo ” deve ser entendido \textbf{aqui} como ponto de partida para iniciar os estudantes no processo de filosofar, o qual não pode ser aplicado como \textbf{um} roteiro de conteúdos fixos a serem transmitidos em sala de aula.
Assim, espera -se \textbf{que} no trabalho com os estudantes, passamos ir para além das orientações presentes neste documento, tanto quanto possível.
Por isso mesmo, propomos enxergar a Filosofia, neste texto, como componente de ensino na educação escolar de nível médio, além de levar os estudantes a lidar com os desafios e as condições de possibilidades \textbf{que} isto \textbf{implica} ; tomar consciência disto é \textbf{imprescindível} para aprender a aprender.
Mas por \textbf{que} estudar filosofia no Ensino Médio?
Entendemos \textbf{que} a Filosofia como componente curricular tem especificidades próprias, distintas da produzida \textbf{enquanto} saber e \textbf{que} \textbf{perder} isto de vista é não entender o papel deste campo do conhecimento neste nível de ensino.
Compreendemos assim, \textbf{que} o essencial é promover \textbf{uma} ação filosófico-educativa constituída pela seleção dos conhecimentos \textbf{que} foram \textbf{historicamente} desenvolvidos no campo da Filosofia de modo a possibilitar ao estudante as condições de se apropriar dos \textbf{fundamentos} \textbf{teóricos} e \textbf{metodológicos} deste saber, assim filosofando em sala de aula. \textbf{Contudo}, não se trata de formar \textbf{um} especialista em Filosofia nesse nível de ensino, e sim, ir aos poucos \textbf{aproximando} os estudantes, em sala de aula, do contato com o estilo reflexivo, problematizador, criativo e argumentativo da Filosofia, \textbf{que} é indispensável para se apreender as formas de \textbf{abordagem} filosófica ( filosofar ) de modo significativo.
Para atender a estes objetivos, a organização curricular \textbf{que} \textbf{ora} apresentamos estrutura -se por eixos temáticos fundamentados na história da filosofia.
Na 1a Série trata -se de organizar as habilidades e competências como \textbf{uma} atitude filosófica para aprender a filosofar ; o foco principal deve ser o de possibilitar aos estudantes \textbf{uma} noção básica e geral daquilo \textbf{que} seja a Filosofia e como se aprende a filosofar.
No 1o bimestre ( Iniciação à Filosofia ) é o momento de apresentar ao estudante a Filosofia ; o modo de fazer isso deve ser definido pelo professor.
No 2o bimestre ( História da Filosofia ) trata -se de iniciar alguma aproximação com aquilo \textbf{que} se entende como a especificidade da filosofia.
Daí sugerir -se fazer isto pela identificação e diferenciação entre o mito e a filosofia como duas formas distintas de explicar a realidade e de orientar o agir humano no mundo.
O 3o bimestre ( Mito, Ruptura e Razão ) é o momento de apresentar ao estudante \textbf{um} dos temas clássicos da filosofia, a questão do \textbf{que} é o \textbf{homem} ( o ser humano ), a partir da qual surgem as primeiras reflexões racionais no campo da ética ( ser singular ) e da política ( ser político ), nasce a antropologia filosófica.
Neste caso, pode -se trabalhar Sócrates e os Sofistas, dentre outros \textbf{autores}.
No 4o bimestre ( Ignorância e Busca da Verdade ) dá -se continuidade à \textbf{abordagem} de outros temas clássicos da filosofia, como a diferença entre doxa e episteme.
Neste momento, parece apropriado trabalhar Platão e a Alegoria da Caverna como explicitação de \textbf{uma} das formas de se entender o processo do filosofar.
A 2a Série, deve continuar com a iniciação ao processo do filosofar com alguns saltos tendo em vista a \textbf{abordagem} de pontos chave necessários à aprendizagem filosófica.
O eixo \textbf{central} desta Série é a questão do aprender a ser e conhecer.
O 1o bimestre ( A compreensão da Lógica ) visa a atender \textbf{uma} necessidade básica da área referente aos modos e meios para se desenvolver o pensar filosófico propriamente \textbf{dito}, para desenvolver o \textbf{raciocínio} lógico e a argumentação.
Para tanto, será necessário recorrer aos princípios e aos instrumentos da lógica formal, bem como é de se esperar \textbf{que} se faça alguma menção à lógica \textbf{dialética}, tanto quanto possível.
O 2o bimestre ( A Perspectiva do Conhecimento ) trata de iniciar com a própria questão do \textbf{que} é o conhecimento, se é possível ao ser humano conhecer e, se isso é possível, em quais condições.
O 3o bimestre ( Da Metafísica à Ciência ) parte da análise e discussão dos vários tipos de conhecimento humano, tais como : senso comum, arte, religião, metafísica, ciência, filosofia etc., com destaque para a relação entre o conhecimento científico e filosófico.
Apesar de ambos pretenderem \textbf{um} conhecimento sistemático e rigoroso da realidade, procura -se \textbf{demarcar} quais são as semelhanças e distinções entre estas duas formas do conhecer humano e os seus respectivos \textbf{métodos}.
O 4o bimestre ( A Atitude Científica ) objetiva fazer a crítica à racionalidade \textbf{moderna} e apontar os seus limites e possibilidades em relação às expectativas criadas no período das Revoluções Científicas, bem como situar as finalidades da escola neste processo tendo em vista ser a escola o “ locus ” por \textbf{excelência} de circulação e de produção do saber \textbf{sistematizado} nos tempos \textbf{atuais}.
Na 3a Série o foco se dirige ao aprender a ser humano e agir, à produção especificamente humana.
Nela se destacaram a Dimensão do Sagrado, a Dimensão Artística, a Dimensão Ética, a Dimensão Política e as relações entre elas.
No 1o bimestre ( Interpretação e \textbf{Contextualização} da Cultura ) trata -se de compreender a cultura como a intervenção na natureza feita pelo \textbf{homem} criando os mais diversos instrumentos e formas de acesso ao mundo, dando sentido a sobrevivência e a existência do indivíduo e da espécie ( logos ).
No 2o bimestre ( A Experiência do Sagrado e o Universo das Artes ) parte -se primeiro da religião como a atividade cultural mais antiga da \textbf{humanidade}, a partir da \textbf{percepção} de \textbf{que} a realidade exterior é algo fora de nosso controle e \textbf{que} a consciência da morte leva à crença em seres e poderes sobrenaturais e, segundo, da necessidade de entender a arte, nas diferentes épocas, como \textbf{uma} das formas do saber humano, para além das Belas Artes.
Chama -se a atenção para a noção ampliada de arte como \textbf{um} bem fazer ( téchné ) \textbf{que} permite a compreensão da vida a partir da arte.
O 3o bimestre ( A Existência Ética ) deve -se partir de \textbf{uma} análise inicial do saber ético \textbf{pontuando} algumas das suas distintas noções para assim discutir os desafios éticos da sociedade \textbf{atual}, diferenciando ética e moral como \textbf{constitutivos} da autonomia e liberdade humanas ( práxis ).
O 4o bimestre ( A Vida Política ) visa a aprofundar a noção de política, das relações de poder e o seu papel na organização da vida social, bem como problematizar sobre a vida centrada em cidades, especialmente no mundo \textbf{contemporâneo} e refletir sobre suas consequências a partir das \textbf{teorias} políticas construídas pela filosofia.
Propõe ainda, \textbf{uma} crítica à ideologia e à alienação presentes nas práticas cotidianas e nas relacionadas ao mundo do trabalho aos Direitos Humanos.
Essa apresentação dos objetos de aprendizagem por série, visa, de certa forma, mais orientar o professor no processo de ensino e aprendizagem do \textbf{que} determinar \textbf{um} caminho. \textbf{Contudo}, a partir das habilidades e competências apresentadas na BNCC, juntamente com a \textbf{maturidade} emocional, física e intelectual do estudante, o destaque dos objetos se torna necessário porque o nível de \textbf{maturidade} de \textbf{uma} estudante ou \textbf{um} estudante do primeiro ano é diferente no terceiro ano, onde podemos falar da cultura, ética, responsabilidade e religião, com \textbf{uma} maior flexibilidade epistêmica em busca de \textbf{uma} fundamentação mais abrangente.
Assim, podemos também, de acordo com a necessidade e a especificidade de cada espaço educativo, abordar a perspectiva do pensamento filosófico a partir de campos e subcampos - bimestrais ou não - \textbf{que} contemplem as várias visões filosóficas, momentos históricos, \textbf{teorias} e concepções de filosofia, bem como os vários pensadores envolvidos naquele tema contemplado.
Na organização dos objetivos de aprendizagem ( \textbf{que} se direcionam aos estudantes ) \textbf{que} segue \textbf{logo} abaixo, dividimos em unidades temáticas e tais unidades contemplam essa forma de \textbf{abordagem} do conhecimento filosófico como \textbf{norteadores} de \textbf{um} caminho, mas apenas \textbf{um} elemento de indicação, por isso mesmo, os objetos de aprendizagem foram preservados em conceitos fundamentais da filosofia.
Desta forma, o professor pode flexibilizar seu fazer pedagógico numa perspectiva mais fluida e livre.
É claro \textbf{que} os temas ficam a critério de cada espaço educativo, \textbf{contudo}, devemos garantir a partir dos direitos de aprendizagem, aquilo \textbf{que} o componente de filosofia propõe \textbf{enquanto} o básico e necessário para se aprender sobre o pensamento filosófico e seus \textbf{desdobramentos} no caminhar histórico-temporal onde a filosofia foi se fazendo e refazendo como primeira interpretação sistemática da realidade em oposição e complementação ao mito, como busca de certezas e verdades prováveis e indubitáveis para o esclarecimento sobre a realidade, a saída do ser humano da menoridade, do senso comum, do misticismo e do pré-conceito.
Nesse sentido, caminhamos pela tradição filosófica ocidental, isso não \textbf{restringe} o professor \textbf{que} desejar incluir no processo de ensino alguns filósofos e filósofas brasileiras ( os ), a filosofia latino-americana, a filosofia africana, etc., resguardando as temáticas filosóficas, os períodos e campos de atuação de cada concepção filosófica dentro da tradição do conhecimento da filosofia \textbf{sistematizada} ocidental e da história do conhecimento sistemático epistemológico.

% matched lemmas: abordagem, aproximar, aqui, atual, autor, central, configurar, constitutivo, contemporâneo, contextualização, contudo, demarcar, desdobramento, despertar, dialético, dizer, enquanto, excelência, fundamento, historicamente, homem, humanidade, implicar, imprescindível, lado, logo, maturidade, metodológico, moderno, método, norteador, ora, percepção, perder, pontuar, preciso, que, raciocínio, restringir, sistematizar, teoria, teórico, um, ver
\end{textsample}
